\chapter{Quantum Algorithms for Linear Differential Equations}
\label{ch:lec13}

\begin{objectives}
\begin{itemize}[leftmargin=1.4em,nosep]
  \item Reduce linear ODEs/PDEs to problems solvable on a quantum computer.
  \item Compare the linear-system, time-marching, and LCHS approaches.
  \item Explain Schr\"odingerization for non-unitary dynamics.
  \item Work the heat equation as a running example.
\end{itemize}
\end{objectives}

\todo{Write the introductory paragraph orienting the lecture.}

\section{Problem setup and semi-discretization}
\label{sec:lec13-1}
\todo{Draft this section.}

\section{Linear-system-of-equations method}
\label{sec:lec13-2}
\todo{Draft this section.}

\section{Time-marching method}
\label{sec:lec13-3}
\todo{Draft this section.}

\section{Linear combination of Hamiltonian simulation (LCHS)}
\label{sec:lec13-4}
\todo{Draft this section.}

\section{Schr\"odingerization}
\label{sec:lec13-5}
\todo{Draft this section.}

\section{Example: the heat equation}
\label{sec:lec13-6}
\todo{Draft this section.}

\exercisesection
\begin{exercise}
\todo{Add exercise.}
\end{exercise}
\begin{exercise}
\todo{Add exercise.}
\end{exercise}
\begin{exercise}
\todo{Add exercise.}
\end{exercise}

\readingsection
\begin{itemize}[leftmargin=1.4em]
  \item Lin Lin, \emph{qasc notes}, Ch.~4.5--4.6.
  \item Lin Lin, \emph{live notes}, Ch.~20.
\end{itemize}
